\documentclass[12pt,a4paper]{article}
\usepackage[utf8]{inputenc}
\usepackage[T1]{fontenc}
\usepackage[french]{babel}
\usepackage{amsmath,amssymb}
\usepackage{geometry}
\geometry{margin=1.8cm}
\usepackage{enumitem}
\usepackage{parskip}
\usepackage{array}
\usepackage{booktabs}
\usepackage{eurosym}
\usepackage{xcolor}
\usepackage{tcolorbox}
\tcbuselibrary{skins,breakable}
\usepackage{tikz}
\usepackage{multicol}
\usepackage{pifont}

% --- Couleurs du jeu ---
\definecolor{bleufonce}{RGB}{20,60,120}
\definecolor{rouge}{RGB}{180,30,30}
\definecolor{vert}{RGB}{20,130,70}
\definecolor{orange}{RGB}{230,120,20}
\definecolor{violet}{RGB}{110,30,140}
\definecolor{grisclair}{RGB}{240,240,240}
\definecolor{or}{RGB}{220,170,20}

% --- Boîtes stylisées ---
\newtcolorbox{regle}[1][]{
  colback=grisclair, colframe=bleufonce, boxrule=1pt,
  arc=3mm, left=3mm, right=3mm, top=2mm, bottom=2mm, #1
}
\newtcolorbox{carte}[2]{
  colback=white, colframe=#1, boxrule=1.2pt,
  arc=2mm, left=3mm, right=3mm, top=2mm, bottom=2mm,
  title=#2, fonttitle=\bfseries
}
\newtcolorbox{reponse}[1][]{
  colback=vert!5, colframe=vert, boxrule=1pt,
  arc=2mm, left=3mm, right=3mm, top=2mm, bottom=2mm, #1
}

\newcommand{\etoiles}[1]{%
  \ifcase#1
  \or \textcolor{or}{\ding{72}}%
  \or \textcolor{or}{\ding{72}\ding{72}}%
  \or \textcolor{or}{\ding{72}\ding{72}\ding{72}}%
  \or \textcolor{or}{\ding{72}\ding{72}\ding{72}\ding{72}}%
  \fi
}

\begin{document}

\begin{center}
{\Huge\bfseries\color{bleufonce} Le Grand Défi des Évolutions}\\[0.3cm]
{\large Jeu mathématique -- Durée : 50 minutes}\\[0.2cm]
{\large Niveau : Première / Terminale -- Taux, coefficients multiplicateurs, taux global et moyen}\\[0.2cm]
{\large\bfseries 22 élèves -- 6 équipes (4 équipes de 4 et 2 équipes de 3)}
\end{center}

\vspace{0.3cm}

% ============================================================
\section*{Règles du jeu}

\begin{regle}
\textbf{But du jeu :} accumuler un maximum de points en résolvant des défis sur les évolutions.\\[0.2cm]
\textbf{Organisation :} la classe est divisée en \textbf{6 équipes} :
\begin{itemize}
    \item Équipes 1, 2, 3 et 4 : 4 élèves chacune (soit 16 élèves).
    \item Équipes 5 et 6 : 3 élèves chacune (soit 6 élèves).
\end{itemize}
Chaque équipe désigne un capitaine qui rend les réponses.\\[0.2cm]
\textbf{Matériel :} une fiche de score par équipe, des cartes-questions, une calculatrice autorisée.\\[0.2cm]
\textbf{Déroulement :} 4 manches chronométrées. À chaque manche, les équipes marquent des points selon le barème indiqué.\\[0.2cm]
\textbf{Victoire :} l'équipe qui totalise le plus de points à la fin des 4 manches remporte le défi.
\end{regle}

\vspace{0.3cm}

\begin{center}
\renewcommand{\arraystretch}{1.4}
\begin{tabular}{|c|l|c|c|}
\hline
\textbf{Manche} & \textbf{Titre} & \textbf{Durée} & \textbf{Points max} \\
\hline
1 & Échauffement éclair & 12 min & 10 pts \\
\hline
2 & Défis en équipe & 15 min & 15 pts \\
\hline
3 & Enveloppe mystère & 12 min & 15 pts \\
\hline
4 & Finale « Tout ou rien » & 11 min & 10 pts \\
\hline
\textbf{Total} & & \textbf{50 min} & \textbf{50 pts} \\
\hline
\end{tabular}
\end{center}

\newpage

% ============================================================
\section*{Manche 1 -- Échauffement éclair (12 min, 10 pts)}

\textit{Chaque bonne réponse rapporte 1 point. Les 10 questions sont posées à toute la classe. La première équipe qui lève la main et donne la bonne réponse marque le point.}\\[0.2cm]
\textbf{Règle d'équité :} chaque équipe ne peut marquer que \textbf{2 points maximum} sur cette manche. Une équipe ayant déjà marqué 2 fois ne peut plus répondre aux questions restantes.

\vspace{0.3cm}

\begin{carte}{bleufonce}{Question 1}
Un prix augmente de \(25\%\). Quel est le coefficient multiplicateur ?
\end{carte}

\begin{carte}{bleufonce}{Question 2}
Un prix baisse de \(10\%\). Quel est le coefficient multiplicateur ?
\end{carte}

\begin{carte}{bleufonce}{Question 3}
Le coefficient multiplicateur est \(1{,}4\). Quel est le taux d'évolution en \% ?
\end{carte}

\begin{carte}{bleufonce}{Question 4}
Le coefficient multiplicateur est \(0{,}75\). Quel est le taux d'évolution en \% ?
\end{carte}

\begin{carte}{bleufonce}{Question 5}
Un article passe de \(50\)~\euro{} à \(60\)~\euro{}. Quel est le taux d'évolution ?
\end{carte}

\begin{carte}{bleufonce}{Question 6}
Une grandeur augmente de \(20\%\), puis de \(20\%\). Le taux global est-il \(40\%\) ?
\end{carte}

\begin{carte}{bleufonce}{Question 7}
Que vaut \(1{,}10 \times 0{,}90\) ?
\end{carte}

\begin{carte}{bleufonce}{Question 8}
Un prix baisse de \(50\%\). Par quel coefficient faut-il multiplier pour revenir au prix initial ?
\end{carte}

\begin{carte}{bleufonce}{Question 9}
Un prix augmente de \(10\%\), puis de \(20\%\). Quel est le coefficient multiplicateur global ?
\end{carte}

\begin{carte}{bleufonce}{Question 10}
Un prix baisse de \(20\%\), puis augmente de \(25\%\). Le prix final est-il égal au prix initial ?
\end{carte}

\newpage

% ============================================================
\section*{Manche 2 -- Défis en équipe (15 min, 15 pts)}

\textit{Chaque équipe reçoit les 3 défis ci-dessous. Elle dispose de 15 minutes pour les résoudre. Chaque défi est noté sur 5 points (3 pts pour la démarche, 2 pts pour le résultat).}

\vspace{0.3cm}

\begin{carte}{rouge}{Défi A}
Un smartphone coûte \(600\)~\euro{}. Son prix baisse de \(15\%\), puis augmente de \(10\%\).\\[0.2cm]
\textbf{Quel est son prix final ? Le prix a-t-il globalement augmenté ou baissé ?}
\end{carte}

\begin{carte}{rouge}{Défi B}
Une ville comptait \(25\,000\) habitants en 2015. Elle a perdu \(4\%\) d'habitants en 2016, puis gagné \(6\%\) en 2017, puis perdu \(2\%\) en 2018.\\[0.2cm]
\textbf{Combien d'habitants compte-t-elle fin 2018 (arrondir à l'unité) ? Quel est le taux global d'évolution ?}
\end{carte}

\begin{carte}{rouge}{Défi C}
Un article subit deux évolutions successives : \(+30\%\) puis \(-30\%\).\\[0.2cm]
\textbf{Le prix final est-il égal au prix initial ? Justifier par un calcul.}
\end{carte}

\newpage

% ============================================================
\section*{Manche 3 -- Enveloppe mystère (12 min, 15 pts)}

\begin{regle}
\textbf{Le défi de la persévérance !}\\[0.2cm]
6 enveloppes (A à F) circulent entre les 6 équipes. Chaque équipe démarre avec une enveloppe différente. Dès qu'elle a terminé, elle la rend et reçoit l'enveloppe suivante (rotation dans le sens horaire).\\[0.2cm]
\textbf{Objectif :} en résoudre le plus possible en 12 minutes.\\[0.3cm]
\textbf{Barème progressif (adapté à 6 équipes) :}
\begin{center}
\renewcommand{\arraystretch}{1.6}
\begin{tabular}{|>{\raggedright\arraybackslash}p{5.5cm}|c|c|}
\hline
\textbf{Enveloppe résolue} & \textbf{Étoiles} & \textbf{Points} \\
\hline
1\textsuperscript{re} enveloppe correcte & \etoiles{1} & \(4\) pts \\
\hline
2\textsuperscript{e} enveloppe correcte & \etoiles{2} & \(5\) pts \\
\hline
3\textsuperscript{e} enveloppe correcte & \etoiles{3} & \(6\) pts \\
\hline
\multicolumn{2}{|r|}{\textbf{Total maximal}} & \(\mathbf{15}\) \textbf{pts} \\
\hline
\end{tabular}
\end{center}
\vspace{0.2cm}
\textbf{Règle d'or :} une enveloppe non résolue ou fausse ne rapporte rien (pas de malus).\\[0.2cm]
\textbf{Astuce :} ne restez pas bloqués ! Si une enveloppe résiste, rendez-la et passez à la suivante.
\end{regle}

\newpage

% --- Enveloppe A (complexifiée) ---
\begin{carte}{violet}{Enveloppe A}
Un commerçant augmente ses prix de \(25\%\). Devant la baisse des ventes, il applique une remise de \(20\%\).\\[0.3cm]
\textbf{1.} Quel est le coefficient multiplicateur global ?\\[0.1cm]
\textbf{2.} Le prix final est-il supérieur, égal ou inférieur au prix initial ? Justifier par un calcul.\\[0.1cm]
\textbf{3.} Quel taux de remise aurait-il fallu appliquer après la hausse de \(25\%\) pour revenir exactement au prix initial ? (arrondir à \(0{,}01\%\))
\end{carte}

\newpage

% --- Enveloppe B (complexifiée) ---
\begin{carte}{violet}{Enveloppe B}
Un article coûte \(360\)~\euro{} après deux évolutions successives : une hausse de \(20\%$, puis une hausse de $t\%$.\\[0.3cm]
\textbf{1.} Déterminer le prix de l'article avant la première hausse.\\[0.1cm]
\textbf{2.} Sachant que le prix initial était de \(250\)~\euro{}, déterminer \(t\) (arrondir à \(0{,}01\%\)).\\[0.1cm]
\textbf{3.} Vérifier le résultat en calculant \(250 \times 1{,}20 \times \left(1+\dfrac{t}{100}\right)\).
\end{carte}

\newpage

% --- Enveloppe C (complexifiée) ---
\begin{carte}{violet}{Enveloppe C}
Un placement a rapporté \(33{,}1\%\) en trois ans, avec le même taux annuel \(t\).\\[0.3cm]
\textbf{1.} Écrire l'équation vérifiée par \(t\).\\[0.1cm]
\textbf{2.} Déterminer \(t\) (arrondir à \(0{,}01\%\)).\\[0.1cm]
\textbf{3.} Un second placement a rapporté \(21\%\) en deux ans, avec le même taux annuel \(t'\). Comparer \(t\) et \(t'\) : quel placement est le plus avantageux ?
\end{carte}

\newpage

% --- Enveloppe D (complexifiée) ---
\begin{carte}{violet}{Enveloppe D}
Une population augmente de \(8\%\) la première année, puis de \(t\%\) la deuxième année. Sur les deux ans, le taux global est de \(20{,}96\%\).\\[0.3cm]
\textbf{1.} Écrire l'équation vérifiée par \(t\).\\[0.1cm]
\textbf{2.} Déterminer \(t\) (arrondir à \(0{,}01\%\)).\\[0.1cm]
\textbf{3.} Calculer le taux annuel moyen sur les deux ans (arrondir à \(0{,}01\%\)) et vérifier qu'il est bien compris entre \(8\%\) et \(t\%\).
\end{carte}

\newpage

% --- Enveloppe E (complexifiée) ---
\begin{carte}{violet}{Enveloppe E}
Une action baisse de \(20\%\), puis remonte de \(x\%\). Après ces deux évolutions, elle est revenue exactement à son cours initial.\\[0.3cm]
\textbf{1.} Justifier que \(0{,}80 \times \left(1+\dfrac{x}{100}\right) = 1\).\\[0.1cm]
\textbf{2.} Déterminer \(x\) (arrondir à \(0{,}01\%\)).\\[0.1cm]
\textbf{3.} Comparer \(x\) avec le taux de baisse initial (\(20\%\)). Que peut-on en conclure de général sur les taux de retour au prix initial ?
\end{carte}

\newpage

% --- Enveloppe F (complexifiée) ---
\begin{carte}{violet}{Enveloppe F}
Deux articles subissent chacun deux évolutions successives.\\[0.2cm]
\begin{itemize}
    \item \textbf{Article 1 :} \(+30\%$, puis $-10\%$.
    \item \textbf{Article 2 :} $+10\%$, puis $+30\%$.
\end{itemize}
\textbf{1.} Calculer le coefficient multiplicateur global de chaque article.\\[0.1cm]
\textbf{2.} Les deux articles ont-ils subi la même évolution globale ? Justifier.\\[0.1cm]
\textbf{3.} Quel taux unique aurait donné le même résultat pour chaque article ? (arrondir à \(0{,}01\%\))
\end{carte}

\newpage

% ============================================================
\section*{Manche 4 -- Finale « Tout ou rien » (11 min, 10 pts)}

\textit{Chaque équipe choisit un niveau de difficulté. Une bonne réponse rapporte les points indiqués ; une mauvaise réponse les retire. Les points peuvent devenir négatifs.}

\vspace{0.3cm}

\begin{center}
\renewcommand{\arraystretch}{1.5}
\begin{tabular}{|c|p{8.5cm}|c|c|}
\hline
\textbf{Niveau} & \textbf{Question} & \textbf{Gain} & \textbf{Perte} \\
\hline
Facile & Un prix augmente de \(30\%\). Quel coefficient appliquer pour revenir au prix initial ? (arrondir à \(0{,}001\)) & \(+2\) & \(-2\) \\
\hline
Moyen & Un capital de \(5\,000\)~\euro{} subit 4 évolutions : \(+8\%\), \(-5\%\), \(+12\%\), \(-10\%\). Quel est le capital final ? & \(+3\) & \(-3\) \\
\hline
Difficile & Une grandeur a subi 5 évolutions annuelles identiques de taux \(t\). Le taux global est \(+27{,}628\%\). Déterminer \(t\). & \(+5\) & \(-5\) \\
\hline
\end{tabular}
\end{center}

\newpage

% ============================================================
\section*{Fiche de score}

\begin{center}
{\large \textbf{Nom de l'équipe :} \dotfill}\\[0.2cm]
{\large \textbf{Capitaine :} \dotfill}
\end{center}

\vspace{0.4cm}

\begin{center}
\renewcommand{\arraystretch}{1.8}
\begin{tabular}{|>{\raggedright\arraybackslash}p{5cm}|c|c|}
\hline
\textbf{Manche} & \textbf{Points max} & \textbf{Points obtenus} \\
\hline
1 -- Échauffement éclair & 10 & \\
\hline
2 -- Défis en équipe & 15 & \\
\hline
3 -- Enveloppe mystère & 15 & \\
\hline
4 -- Finale « Tout ou rien » & 10 & \\
\hline
\textbf{Total} & \textbf{50} & \\
\hline
\end{tabular}
\end{center}

\vspace{0.5cm}

\section*{Détail Manche 3 -- Enveloppe mystère}

\begin{center}
\renewcommand{\arraystretch}{1.6}
\begin{tabular}{|c|c|c|c|}
\hline
\textbf{Ordre} & \textbf{Enveloppe} & \textbf{Correcte ? (O/N)} & \textbf{Points} \\
\hline
1\textsuperscript{re} \etoiles{1} & A \quad B \quad C \quad D \quad E \quad F & & /4 \\
\hline
2\textsuperscript{e} \etoiles{2} & A \quad B \quad C \quad D \quad E \quad F & & /5 \\
\hline
3\textsuperscript{e} \etoiles{3} & A \quad B \quad C \quad D \quad E \quad F & & /6 \\
\hline
\multicolumn{3}{|r|}{\textbf{Total Manche 3}} & \textbf{/15} \\
\hline
\end{tabular}
\end{center}

\vspace{0.5cm}

\begin{regle}
\textbf{Conseils pour l'enseignant :}
\begin{itemize}
    \item Distribuer la fiche de score dès le début.
    \item Chronométrer chaque manche rigoureusement pour respecter les 50 minutes.
    \item Pour la manche 1, utiliser un buzzer ou un simple lever de main. Rappeler la règle des 2 points maximum par équipe.
    \item Pour la manche 3, prévoir au moins 2 copies de chaque enveloppe pour fluidifier la rotation (12 copies au total).
    \item Pour la manche 4, annoncer les points avant la réponse pour créer du suspense.
    \item En cas d'égalité, poser la question bonus : « Un prix augmente de \(100\%\), puis baisse de \(50\%\). Revient-il au prix initial ? »
\end{itemize}
\end{regle}

\vspace{0.5cm}

\begin{center}
{\Large\bfseries\color{bleufonce} Bon jeu aux 6 équipes !}
\end{center}

\end{document}